%% kapa: exam=2 round=24 year=2013 period=1 subject=appraisal_practice subject_ko=감정평가실무 seq=1 number=1 points=35
%% source: Q-Net 감정평가사 기출문제 (cst00309) artlSeq=5077404 file=extracted/1008/5077404/2013년 제24회 감정평가사 2차시험 기출문제/1교시.hwp
%% conversion: kordoc 4.24.51 -> tools/gen2.py (auto); review=proofread-sonnet; flags=none
\begin{longquestion}{【문제 1】}
토지 소유자인 甲 법인은 골프장 개발업체인 乙 법인과 다음과 같은 계약을 맺었다. \par
\begin{framedtext}
<계약내용> \par
ㆍ乙 법인은 甲 법인의 토지를 임차하여 골프장(27홀)으로 개발하여 운영한다. \par
ㆍ골프장 개발과 관련된 인허가 비용은 甲 법인 부담으로 하고 개발비용은 乙 법인 부담으로 한다. \par
ㆍ乙 법인은 골프장 준공일로부터 연간 토지임대료 1,000,000,000원을 甲 법인에 매년 초 지급하며 연간 2\%씩 임대료를 상승하여 지급한다. \par
ㆍ골프장 운영과 관련된 제반 유지보수비용, 보험료, 제세공과 등은 운영자 부담으로 한다. \par
ㆍ계약기간은 준공일로부터 10년이고, 계약기간 만료일 乙 법인이 개발한 모든 골프장 시설 등은 甲 법인으로 귀속되며, 甲 법인은 乙 법인의 최초 개발비용의 30\% 상당액을 乙 법인에 지급한다. \par
\end{framedtext}
乙 법인의 골프장 개발계획은 순조롭게 진행되어 2013. 1. 1에 준공하였다. 다음 물음에 답하시오. (35점) \par
\medskip\noindent 물음 1) 감정평가사 丙 씨는 甲 법인으로부터 2013. 1. 1자 甲 법인 소유토지에 대한 가치산정을 의뢰 받았다. 주어진 자료를 활용하여 가치를 산정하고 평가방법에 대해 서술하시오. (20점) \par
\medskip\noindent 물음 2) 이러한 계약을 하는 甲 법인과 乙 법인은 합리적 의사결정을 하는 것인가에 대해 NPV 법으로 검토하고 서술하시오. (15점) \par
\textbf{<자료 1> 토지목록(甲 법인)} \par
{\small\centering
\begin{tabular}{|c|c|c|c|c|c|}
\hline
\makecell{기호} & \makecell{소재지} & \makecell{지번} & \makecell{지목} & \makecell{면적(m\textsuperscript{2})} & \makecell{용도지역} \\
\hline
\makecell{1} & \makecell{A군 B면 C리} & \makecell{200} & \makecell{전} & \makecell{1,000} & \makecell{계획관리} \\
\hline
\makecell{2} & \makecell{A군 B면 C리} & \makecell{200-1} & \makecell{전} & \makecell{2,600} & \makecell{계획관리} \\
\hline
\makecell{3} & \makecell{A군 B면 C리} & \makecell{200-2} & \makecell{전} & \makecell{1,550} & \makecell{계획관리} \\
\hline
\makecell{4} & \makecell{A군 B면 C리} & \makecell{200-3} & \makecell{답} & \makecell{2,350} & \makecell{계획관리} \\
\hline
\makecell{5} & \makecell{A군 B면 C리} & \makecell{200-5} & \makecell{전} & \makecell{1,300} & \makecell{계획관리} \\
\hline
\makecell{6} & \makecell{A군 B면 C리} & \makecell{200-6} & \makecell{전} & \makecell{1,600} & \makecell{계획관리} \\
\hline
\makecell{7} & \makecell{A군 B면 C리} & \makecell{201} & \makecell{전} & \makecell{1,750} & \makecell{계획관리} \\
\hline
\makecell{8} & \makecell{A군 B면 C리} & \makecell{202} & \makecell{답} & \makecell{3,700} & \makecell{계획관리} \\
\hline
\makecell{9} & \makecell{A군 B면 C리} & \makecell{산 100-1} & \makecell{임야} & \makecell{4,500} & \makecell{계획관리} \\
\hline
\makecell{10} & \makecell{A군 B면 C리} & \makecell{산 100-2} & \makecell{임야} & \makecell{1,500,000} & \makecell{계획관리} \\
\hline
\makecell{11} & \makecell{A군 B면 C리} & \makecell{산 100-3} & \makecell{전} & \makecell{900} & \makecell{계획관리} \\
\hline
\makecell{계} &  &  &  & \makecell{1,521,250} &  \\
\hline
\end{tabular}\par}

※ 사업승인면적은 1,450,000m\textsuperscript{2}이며, 나머지는 산 100-2번지 일부로서 자연림 상태의 원형을 유지하고 있음 \par
\textbf{<자료 2> 표준지 공시지가} \par
(2013. 1. 1) \par
{\small\centering
\setlength{\tabcolsep}{2.5pt}
\begin{tabular}{|c|c|c|c|c|c|c|c|c|c|}
\hline
\makecell{일련번호} & \makecell{소재지} & \makecell{지번} & \makecell{면적\newline{}(m\textsuperscript{2})} & \makecell{지목} & \makecell{이용\newline{}상황} & \makecell{용도\newline{}지역} & \makecell{도로\newline{}교통} & \makecell{형상\newline{}지세} & \parbox{1.7cm}{\centering 공시지가(원/m\textsuperscript{2)}} \\
\hline
\makecell{1} & \makecell{A군 D면 E리} & \makecell{14\newline{}(××골프장)} & \makecell{201,000\newline{}(일단지)} & \makecell{임야} & \makecell{골프장} & \makecell{계획\newline{}관리} & \makecell{소로\newline{}한면} & \makecell{부정형\newline{}완경사} & \makecell{50,000} \\
\hline
\makecell{2} & \makecell{A군 B면 C리} & \makecell{190} & \makecell{4,627} & \makecell{전} & \makecell{전} & \makecell{계획\newline{}관리} & \makecell{세로\newline{}(가)} & \makecell{부정형\newline{}완경사} & \makecell{30,000} \\
\hline
\makecell{3} & \makecell{A군 B면 C리} & \makecell{210} & \makecell{1,096} & \makecell{답} & \makecell{답} & \makecell{계획\newline{}관리} & \makecell{세로\newline{}(불)} & \makecell{부정형\newline{}완경사} & \makecell{20,000} \\
\hline
\makecell{4} & \makecell{A군 B면 C리} & \makecell{산 110} & \makecell{32,000} & \makecell{임야} & \makecell{임야} & \makecell{계획\newline{}관리} & \makecell{세로\newline{}(불)} & \makecell{부정형\newline{}완경사} & \makecell{12,000} \\
\hline
\end{tabular}\par}

※ 상기 표준지 공시지가는 2013. 1. 1 고시된 것으로 봄 \par
\textbf{<자료 3> 요인비교} \par
1. 본건 준공된 골프장과 일련번호 1 비교표준지는 제반여건이 유사하나 접근성 등에서 본건이 약 2\% 우세함 \par
2. 소지상태로서의 본건은 일련번호 2～4와 비교시 지목별로 대체로 유사하나 본건 “전”은 표준지 대비 3\%, 본건 “답”은 표준지 대비 2\%, 본건 “임야”는 표준지 대비 1\%씩 각각 우세함 \par
3. 주변 거래사례나 평가선례를 분석하여 보면 표준지 일련번호 1은 기타요인보정이 필요 없으나 표준지 일련번호 2, 3은 10\%, 표준지 일련번호 4는 20\% 기타요인보정이 필요함 \par
\textbf{<자료 4> 제 비용} \par
1. 인허가 관련비용: \par
인허가 비용 1,500,000,000원 \par
제반부담금 2,500,000,000원 \par
제세공과 500,000,000원 \par
기타 3,000,000,000원 \par
계 7,500,000,000원 \par
2. 골프장 조성(개발)공사 비용: 홀당 1,400,000,000원 \par
\textbf{<자료 5> 주변사례 및 시장동향} \par
주변 유사 골프장의 사례를 보면 연간 9홀 기준으로 20억원의 영업이익이 발생하는 것으로 조사되었으며, 향후 매년 영업이익증가율은 1\% 정도 일 것으로 추정됨. 그러나 시장의 수요ㆍ공급을 예측하여 보면 연간 9홀 기준으로 22억원이 한계점인 것으로 조사됨 \par
\textbf{<자료 6> 기타} \par
1. 클럽하우스 등 건물은 고려하지 아니함 \par
2. 인허가 및 개발에 관한 제 비용은 준공일에 투입된 것으로 가정하며, 영업이익의 발생시점은 토지임대료를 지급하는 매년 초에 발생하는 것으로 가정함 \par
3. 준공 후 제반세금은 고려하지 아니하며 영업이익은 현금유입액으로 봄 \par
4. 본건 투자에 있어 타인자본은 고려하지 아니함 \par
5. 현금흐름(cash flow)에 적용된 할인율은 연 7\%로 조사되었음 \par
6. 기말복귀액 산정시 영업이익과 골프장 가치에 적용할 환원이율은 연 8\%로 조사되었음 \par
7. 본건 주변 표준지 공시지가의 향후 10년 후 예상상승률은 10\% 정도임 \par
\textbf{<자료 7> 시간가치율} \par
1. 현재 가치율(7\%) \par
{\small\centering
\begin{tabular}{|c|c|c|c|c|c|c|c|c|c|c|}
\hline
\makecell{1기초} & \makecell{2기초} & \makecell{3기초} & \makecell{4기초} & \makecell{5기초} & \makecell{6기초} & \makecell{7기초} & \makecell{8기초} & \makecell{9기초} & \makecell{10기초} & \makecell{10기말} \\
\hline
\makecell{1.000} & \makecell{0.9346} & \makecell{0.8734} & \makecell{0.8163} & \makecell{0.7629} & \makecell{0.7130} & \makecell{0.6663} & \makecell{0.6227} & \makecell{0.5820} & \makecell{0.5439} & \makecell{0.5083} \\
\hline
\end{tabular}\par}

2. 미래 가치율(2\%) \par
{\small\centering
\begin{tabular}{|c|c|c|c|c|c|c|c|c|c|c|}
\hline
\makecell{1.02\textsuperscript{0}} & \makecell{1.02\textsuperscript{1}} & \makecell{1.02\textsuperscript{2}} & \makecell{1.02\textsuperscript{3}} & \makecell{1.02\textsuperscript{4}} & \makecell{1.02\textsuperscript{5}} & \makecell{1.02\textsuperscript{6}} & \makecell{1.02\textsuperscript{7}} & \makecell{1.02\textsuperscript{8}} & \makecell{1.02\textsuperscript{9}} & \makecell{1.02\textsuperscript{10}} \\
\hline
\makecell{1.00} & \makecell{1.0200} & \makecell{1.0404} & \makecell{1.0612} & \makecell{1.0824} & \makecell{1.1041} & \makecell{1.1262} & \makecell{1.1487} & \makecell{1.1717} & \makecell{1.1951} & \makecell{1.2190} \\
\hline
\end{tabular}\par}

3. 미래 가치율(1\%) \par
{\small\centering
\begin{tabular}{|c|c|c|c|c|c|c|c|c|c|c|}
\hline
\makecell{1.01\textsuperscript{0}} & \makecell{1.01\textsuperscript{1}} & \makecell{1.01\textsuperscript{2}} & \makecell{1.01\textsuperscript{3}} & \makecell{1.01\textsuperscript{4}} & \makecell{1.01\textsuperscript{5}} & \makecell{1.01\textsuperscript{6}} & \makecell{1.01\textsuperscript{7}} & \makecell{1.01\textsuperscript{8}} & \makecell{1.01\textsuperscript{9}} & \makecell{1.01\textsuperscript{10}} \\
\hline
\makecell{1.00} & \makecell{1.0100} & \makecell{1.0201} & \makecell{1.0303} & \makecell{1.0406} & \makecell{1.0510} & \makecell{1.0615} & \makecell{1.0721} & \makecell{1.0829} & \makecell{1.0937} & \makecell{1.1046} \\
\hline
\end{tabular}\par}

\end{longquestion}
